We compared testing rates between White and Black infants by year and within sub-populations defined by maternal Substance Use Disorder (SUD) history and the county prevalence of neonatal abstinence syndrome. We also examined positive test rates among tested infants by racial group and sub-population, which is one way to assess concerns that different clinical standards are used to decide whether to order a test for White and Black infants.  

We use an inverse propensity score weighting strategy to decompose racial differences in testing rates into a component that reflects group differences in risk factors, and a component that may reflect structural differences in the way that health care providers respond to the same observed risk factors differently in White and Black infants \citep{fortin2011decomposition}. 

To fix ideas, suppose $X_i$ represents the set of clinically relevant factors that are observed by the clinician and are also contained in the EMR. $U_i$ represents additional factors that may be observed by the clinician but are not available in the data. $T_i$ is a binary variable indicating whether the clinician orders the test. To distinguish between racial group differences in testing rates that are due to clinical risks rather than differential treatment, assume test decisions are determined by $T_i = V_{r}(X_i, U_i)$, where $V_{r}(\cdot)$ is the physician's testing decision function -- which could differ by race, $r =[W,B]$. 

Under the assumption that unobserved clinical factors are independent of the infant's recorded race conditional on covariates, aggregate differences in testing rates across White and Black infants may arise for two reasons: (i) systematic differences in clinically relevant traits represented by $X_i$ and $U_i$, and (ii) structural differences in the way that clinicians interpret clinically relevant traits for White and Black infants. In this notation, the overall difference in testing rates between White and Black infants is $\Delta_{B,W} = E[V_B(x,u)|B = 1] - E[V_W(x,u)|W = 1]$. This gap in testing rates can be decomposed into two pieces:

\begin{align*}
    \Delta_{B,W} &= \underbrace{E[V_B(x,u) - V_W(x,u) \mid W_i = 1]}_{\delta_s} + \underbrace{\big(E[V_B(x,u) \mid B_i = 1] - E[V_B(x,u) \mid W_i = 1]\big)}_{\delta_x}
\end{align*}

In the formula, $\delta_s$ represents the difference in the testing rates that would prevail if Black infants had the same observable risk factors as White infants, but these risks were evaluated using race-specific decision functions. Conversely, $\delta_x$ represents the difference in testing rates that would be observed if White and Black infants possessed their actual observed risk factors but were both assessed using the Black decision function. Conceptually, $\delta_x$ is the part of the testing rate gap that arises because White and Black infants have different clinical risk factors. $\delta_s$ is the part of the observed racial testing gap that is not explained by observable differences in risk factors. It arises because clinicians interpret and respond to the same clinical information in different ways for White and Black infants, or because clinicians observe additional risk factors that are not recorded in the data. The share of the overall gap that is explained by risk factors is $\frac{\delta_x}{\Delta_{B,W}}$. The unexplained share is $\frac{\delta_s}{\Delta_{B,W}}$. We use an inverse propensity score weighting method to perform these decompositions and to estimate the structural and clinical component of racial testing gaps.\footnote{In the appendix, we also report estimates from Oaxaca-Blinder decompositions, which produces a similar decomposition using parametric regressions.} 
